\documentclass[11pt]{article}
\usepackage[margin=1in]{geometry}
\usepackage[T1]{fontenc}
\usepackage[utf8]{inputenc}
\usepackage{hyperref}
\usepackage{longtable}
\usepackage{booktabs}
\usepackage{xcolor}
\usepackage{array}
\usepackage{enumitem}

\setlength{\parindent}{0pt}
\setlength{\parskip}{0.7em}

\definecolor{passgreen}{RGB}{0,128,0}
\definecolor{failred}{RGB}{180,0,0}
\definecolor{warnamber}{RGB}{180,100,0}

\newcommand{\pass}{\textcolor{passgreen}{\textbf{PASS}}}
\newcommand{\fail}{\textcolor{failred}{\textbf{CRITICAL FAIL}}}
\newcommand{\warn}{\textcolor{warnamber}{\textbf{WARN}}}
\newcommand{\sanfail}{\textcolor{failred}{\textbf{FAIL (sanity check)}}}

\title{Implementation Audit Report\\Long Island Traffic Accident Cluster Discovery System\\
\large Current State, Verified Results, and Required Fixes}
\author{AMS 597 Project Team}
\date{\today}

\begin{document}
\maketitle
\tableofcontents
\newpage

% ============================================================
\section{Purpose of This Document}
% ============================================================

This document is a post-execution audit of the Long Island accident clustering workflow
implemented in \texttt{li\_clustering.Rmd}. It documents the current state of the
implementation after a full run under the \texttt{local\_safe} profile, evaluates each
phase against the original Software Requirements Specification
(\texttt{Li\_Clustering\_SRS.tex}), records all verified tangible outputs, identifies
failures and deficiencies, and specifies in precise terms what must be fixed before
the work can be considered academically defensible.

This is not a replacement for the SRS. The SRS defines what the system \textit{should do}.
This audit documents what it \textit{actually does} and where the gap lies.

% ============================================================
\section{Workflow Architecture}
% ============================================================

The implementation is structured as a ten-phase reproducible R Markdown document. The
top-level pipeline is:

\begin{center}
Environment setup
$\rightarrow$ Data ingestion and profiling
$\rightarrow$ Feature engineering
$\rightarrow$ Matrix construction
$\rightarrow$ k-Prototypes candidate runs
$\rightarrow$ Cluster validation and stability
$\rightarrow$ DBSCAN spatial clustering
$\rightarrow$ Post-clustering interpretation
$\rightarrow$ Export and visualization
$\rightarrow$ Submission checklist
\end{center}

Two independent clustering tracks are maintained throughout:

\begin{itemize}
\item \textbf{Scenario clustering}: \texttt{kproto(type = "gower")} on 16 mixed-type
  contextual variables (time, weather, roadway context, environmental conditions)
\item \textbf{Spatial clustering}: \texttt{dbscan()} on projected kilometer-coordinate
  representations of accident start locations
\end{itemize}

The two tracks are merged at Phase 8 for cross-tabulation and joint interpretation.

% ============================================================
\section{Phase-by-Phase Audit Summary}
% ============================================================

\begin{longtable}{|p{3.2cm}|p{2cm}|p{8.5cm}|}
\hline
\textbf{Phase} & \textbf{Status} & \textbf{Summary} \\
\hline
\endfirsthead
\hline
\textbf{Phase} & \textbf{Status} & \textbf{Summary} \\
\hline
\endhead
Phase 1 --- Environment Setup & \pass & Reproducibility solid. Seed fixed at 42.
Run profiles defined correctly. \texttt{local\_safe} active. All package checks
pass. \\ \hline
Phase 2 --- Data Ingestion & \pass & Dimension assertions (33{,}375 rows
$\times$ 48 columns) hardcoded and verified. Missingness audit complete. Severity
distribution confirmed. Binary flag prevalence profiled. \\ \hline
Phase 3 --- Feature Engineering & \pass & Temporal, weather, wind, and binary
roadway features engineered correctly. Hour-band, season, rush-hour, and
weather-group derivations are methodologically sound. \\ \hline
Phase 4 --- Matrix Construction & \pass & Mixed-type scenario matrix built with
explicit type enforcement and post-imputation NA check. Spatial km-projection
formulas are geometrically correct for the Long Island region. \\ \hline
Phase 5 --- k-Prototypes Candidates & \pass & Candidate solutions for
$k = 3, 4, 5, 6, 7, 8$ generated with 10 random starts each. Within-cluster
distance diagnostics and dominance flags saved. \\ \hline
Phase 6 --- Cluster Validation & \fail & Silhouette scores 0.15--0.18 (below the
0.30 threshold). Jaccard stability 0.20 (below the 0.60 threshold). k-selection
fallback logic has a suspected bug. No explicit user-facing warning on quality
failure is raised. \\ \hline
Phase 7 --- DBSCAN Spatial Clustering & \pass & 12 corridor clusters produced.
Noise 15.2\%. Largest cluster 44.3\%. All three SRS-defined acceptance criteria
met. Parameter grid search systematic and justified. \\ \hline
Phase 8 --- Post-Clustering Interpretation & \pass & Severity profiling,
county composition, weather and time profiles, and spatial-scenario
cross-tabulation all complete. High-risk clusters 4 and 5 correctly identified
relative to baseline severity rates. \\ \hline
Phase 9 --- Export and Visualization & \pass & Full figure suite, CSV outputs,
and interactive HTML widgets generated. Cluster-labeled dataset
\texttt{li\_accidents\_clustered.csv} saved at 33{,}375 rows $\times$ 58 columns. \\ \hline
Phase 10 --- Submission Checklist & \sanfail & Sanity check 4 of 7
(Jaccard $\geq$ 0.4) explicitly fails with observed value 0.1989. Limitations
table in the same phase does not name the cluster-quality failure explicitly. \\
\hline
\end{longtable}

% ============================================================
\section{Detailed Phase Findings}
% ============================================================

\subsection{Phase 1: Environment and Reproducibility Setup}

\textbf{Implementation lines:} 93--232.

\textbf{What it does:} Specifies required packages, installs missing ones (or halts),
sets global seed, defines run profiles, creates output directories, and validates
input file existence.

\textbf{Verified behavior:}
\begin{itemize}
\item Seed \texttt{set.seed(42)} applied globally at line 130.
\item Three run profiles defined: \texttt{local\_safe}, \texttt{consensus\_16gb},
  \texttt{full\_srs}.
\item Profile controls Phase 6 sample sizes: 5{,}000 rows (local\_safe),
  4{,}000 rows $\times$ 4 repeats (consensus\_16gb), full 33{,}375 rows
  (full\_srs).
\item Package installation check halts clearly on missing packages.
\end{itemize}

\textbf{Issues:}
\begin{itemize}
\item The active profile is \texttt{local\_safe} (line 135). This uses only 15\% of the
  data for validation, which is insufficient for final results. The SRS recommends
  \texttt{consensus\_16gb} as the minimum defensible profile for final submission.
\item No memory check is performed before the run. On machines with limited RAM,
  the \texttt{full\_srs} profile will exceed the 24 GB R vector limit
  (observed in testing).
\end{itemize}

\subsection{Phase 2: Data Ingestion and Raw Profiling}

\textbf{Implementation lines:} 247--430.

\textbf{What it does:} Loads cleaned dataset, asserts exact dimensions, audits
county and severity distributions, profiles missingness and near-zero binary
flags, and inspects duration and weather distributions.

\textbf{Verified behavior:}
\begin{itemize}
\item Dimension assertions: 33{,}375 rows $\times$ 48 columns with \texttt{stop()}
  on mismatch.
\item County counts: Nassau 19{,}390, Suffolk 13{,}838, Queens 147.
\item Severity: Sev1=161, Sev2=28{,}357, Sev3=3{,}938, Sev4=919.
\item Missingness: \texttt{End\_Lat}/\texttt{End\_Lng} 39.2\%, \texttt{Precipitation(in)}
  20.1\%, \texttt{Wind\_Chill(F)} 15.0\%.
\item Zero-prevalence flags (\texttt{Bump}, \texttt{Roundabout}, \texttt{Turning\_Loop},
  \texttt{Traffic\_Calming}, \texttt{No\_Exit}) correctly identified for exclusion.
\end{itemize}

\textbf{Issues:} None. Phase 2 is complete and defensible.

\subsection{Phase 3: Feature Engineering}

\textbf{Implementation lines:} 432--618.

\textbf{What it does:} Derives temporal features from \texttt{Start\_Time}, applies
99th-percentile cap to \texttt{Duration\_Minutes} then log-transforms it, normalizes
weather condition strings into seven groups, standardizes wind direction to nine
compass levels, and identifies which binary roadway flags have sufficient prevalence
to include in clustering.

\textbf{Verified behavior:}
\begin{itemize}
\item Hour-band encoding: 6 ordered levels (Early\_AM, AM\_Rush, Midday, PM\_Rush,
  Evening, Late\_Night).
\item Season: 4 levels (Winter, Spring, Summer, Fall) derived from month.
\item Rush-hour: TRUE for hours 6--9 and 15--18.
\item Duration: capped at 99th percentile (approximately 1{,}445 minutes), then
  \texttt{log1p()}-transformed. Median log value $\approx$ 3.29.
\item Weather groups: Clear, Cloudy, Rain, Snow\_Ice, Fog\_Haze, Thunderstorm, Other.
  Pattern matching via \texttt{grepl()} for Thunderstorm; all others via explicit
  lookup.
\item Wind direction: 9 compass levels plus Variable\_Calm. Both abbreviated
  (\texttt{"N"}) and full-name (\texttt{"North"}) inputs handled.
\item Retained binary flags: \texttt{Crossing}, \texttt{Junction},
  \texttt{Traffic\_Signal}, \texttt{Stop}. All others excluded.
\end{itemize}

\textbf{Issues:}
\begin{itemize}
\item Lines 585--589: Retained binary flags are converted to factor using
  \texttt{factor(..., labels = c("FALSE", "TRUE"))} rather than the simpler
  \texttt{as.factor()}. This is unconventional and can confuse downstream code
  that inspects factor labels. It does not cause a functional error in the current
  run but should be corrected.
\item Thunderstorm matching relies on raw data containing the literal string
  ``Thunderstorm.'' If weather condition strings are abbreviated differently in
  a new dataset, Thunderstorm entries will fall into the \texttt{Other} group
  silently.
\end{itemize}

\subsection{Phase 4: Build Clustering Input Matrices}

\textbf{Implementation lines:} 620--812.

\textbf{What it does:} Assembles the final scenario feature matrix with 11 factor
columns and 5--6 numeric columns, imputes missing values, validates that no NAs
remain, and projects geographic coordinates into local km units for DBSCAN.

\textbf{Verified behavior:}
\begin{itemize}
\item Factor columns (11): \texttt{hour\_band}, \texttt{weekday}, \texttt{weekend},
  \texttt{season}, \texttt{rush\_hour}, \texttt{Sunrise\_Sunset},
  \texttt{weather\_group}, \texttt{Crossing}, \texttt{Junction},
  \texttt{Traffic\_Signal}, \texttt{Stop}.
\item Numeric columns (6): \texttt{Temperature(F)}, \texttt{Humidity(\%)},
  \texttt{Pressure(in)}, \texttt{Visibility(mi)}, \texttt{Wind\_Speed(mph)},
  \texttt{log\_Duration}.
\item Imputation: median for numeric columns, mode for factor columns with fewer
  than 10 missing values, new \texttt{Unknown} level for factor columns with
  10 or more missing values.
\item Post-imputation assertion: total NAs must equal zero, halts otherwise.
\item Coordinate projection reference: $(40.8^\circ\text{N},\; 73.2^\circ\text{W})$.
  Formulas: $x_\text{km} = (\text{lng} - \text{lng}_0) \times 111.32 \times
  \cos(\text{lat}_0 \times \pi / 180)$, $y_\text{km} = (\text{lat} - \text{lat}_0)
  \times 110.574$. These are correct for local-scale projections on Long Island.
\item Projected span: approximately 70 km east-west, 38 km north-south. Consistent
  with Long Island geography.
\end{itemize}

\textbf{Issues:}
\begin{itemize}
\item The imputation threshold (10 missing values) is arbitrary and not documented
  in the SRS. A brief justification note should be added to the report.
\item The coordinate reference point is not documented. A comment explaining
  why $(40.8^\circ\text{N},\; 73.2^\circ\text{W})$ was chosen would improve
  reproducibility.
\end{itemize}

\subsection{Phase 5: k-Prototypes Scenario Clustering}

\textbf{Implementation lines:} 814--971.

\textbf{What it does:} Runs \texttt{kproto()} with Gower distance for
$k = 3$ through $8$, uses 10 random starts, stores within-cluster distance
totals and cluster size diagnostics, flags potentially unstable solutions
(any cluster smaller than 200 rows), and saves an elbow plot.

\textbf{Verified behavior:}

\begin{center}
\begin{tabular}{|c|r|r|r|l|}
\hline
$k$ & tot.withinss & min cluster & max cluster & Unstable \\
\hline
3 & 85{,}386 & 9{,}312 & 13{,}470 & No \\
4 & 80{,}678 & 5{,}126 & 10{,}584 & No \\
5 & 78{,}654 & 2{,}319 & 10{,}509 & No \\
6 & 75{,}187 & 2{,}562 & 8{,}994  & No \\
7 & 72{,}539 & 2{,}089 & 9{,}126  & No \\
8 & 69{,}535 & 2{,}137 & 8{,}543  & No \\
\hline
\end{tabular}
\end{center}

\textbf{Issues:}
\begin{itemize}
\item The within-cluster distance decreases monotonically across $k = 3$--$8$.
  No visible elbow is present, making elbow-based selection ambiguous.
\item At $k = 5$ the min/max cluster ratio is approximately 1:4.5. The solution
  is not flagged as unstable (above 200 threshold) but the imbalance is notable
  and not highlighted.
\item No variance-explained ($R^2$) metric is computed in this phase. This would
  be a useful complement to the within-cluster distance plot.
\end{itemize}

\subsection{Phase 6: Cluster Validation and Stability}
\label{sec:phase6}

\textbf{Implementation lines:} 973--1390.

\textbf{What it does:} Runs \texttt{validation\_kproto()} on a sample of the data
to compute silhouette widths for each candidate $k$, selects the best $k$ by
silhouette, runs \texttt{stability\_kproto()} via bootstrap resampling to compute
Rand and Jaccard indices, and applies a stability-based fallback selection rule.

\subsubsection{Observed Validation Results}

\begin{center}
\begin{tabular}{|c|r|}
\hline
$k$ & Silhouette width \\
\hline
3 & 0.1580 \\
4 & 0.1492 \\
5 & 0.1769 \\
6 & 0.1514 \\
7 & 0.1615 \\
8 & 0.1532 \\
\hline
\end{tabular}
\end{center}

\textbf{Assessment:} All silhouette scores fall below the standard interpretive
threshold of 0.30 (Kaufman and Rousseeuw, 1990). A silhouette of 0.15--0.18 indicates
that cluster boundaries are weak and that many points are nearly equidistant from
their assigned cluster and the nearest alternative cluster. This is not a passing
result by published standards.

\subsubsection{Observed Stability Results}

\begin{center}
\begin{tabular}{|c|r|r|}
\hline
$k$ & Rand index & Jaccard index \\
\hline
5 & 0.7104 & 0.1989 \\
7 & 0.7794 & 0.1666 \\
\hline
\end{tabular}
\end{center}

\textbf{Assessment:} The Rand index values are moderate (0.71, 0.78). However,
the Rand index is known to be inflated for large datasets with many clusters
because most pairs of points are correctly labeled as being in different clusters
by chance. The Jaccard index is a more conservative and informative metric here.
Jaccard values below 0.2 indicate that the cluster assignments change substantially
across bootstrap samples. Published thresholds are: Jaccard $\geq 0.75$ for
stable, $\geq 0.60$ for acceptable, $< 0.60$ for weak (Hennig, 2007). Both
observed values (0.1989 and 0.1666) are far below the acceptable threshold.

\subsubsection{k-Selection Logic Bug}

The fallback selection rule at approximately line 1288 reads:

\begin{verbatim}
if (length(selected_stability) == 1 &&
    selected_stability < 0.5 &&
    nrow(stability_summary) > 1) {
  selected_k <- stability_summary[order(-jaccard, -rand, k)][1, k]
}
\end{verbatim}

All three conditions should be true when running on the observed results:
\texttt{length(selected\_stability) == 1} (exactly one value extracted),
\texttt{selected\_stability = 0.1989 < 0.5} (true), and
\texttt{nrow(stability\_summary) = 2 > 1} (true). Therefore the fallback
\textit{should} override the silhouette selection (which chose $k=5$) and
select $k=7$ based on higher Jaccard. However, the output confirms $k=5$
was used for Phase 8 profiles.

Possible explanations:
\begin{itemize}
\item \texttt{selected\_stability} is extracted as \texttt{NA} rather than 0.1989
  due to a data.table key mismatch, causing \texttt{length()} to return 0 or
  the comparison to produce \texttt{NA} rather than \texttt{TRUE}.
\item The condition evaluates in an unexpected order due to R's short-circuit
  evaluation and NA propagation.
\end{itemize}

\textbf{This is a defect that must be diagnosed and fixed.}

\subsubsection{Missing Quality Gate}

The current implementation selects $k=5$ despite failing both primary quality
thresholds (silhouette and Jaccard). No warning is printed, no conditional halt
or explicit exploratory flag is set, and the subsequent phases proceed as if
the solution were adequate. The only acknowledgment is a note in the Phase 10
sanity check comment text, which a reader of the report would not encounter.

\subsection{Phase 7: DBSCAN Spatial Hotspot Clustering}

\textbf{Implementation lines:} 1392--1721.

\textbf{What it does:} Constructs a systematic grid search over
$\texttt{eps} \in [0.3, 2.0]$ km (step 0.1) and
$\texttt{minPts} \in \{40, 50, 75, 100\}$, evaluates each combination
against three SRS acceptance criteria, selects the combination closest to
the target operational point, and runs the final DBSCAN solution.

\textbf{Verified results:}
\begin{itemize}
\item Selected parameters: \texttt{eps} = 1.0 km, \texttt{minPts} = 100.
\item Number of non-noise clusters: 12 (SRS target: 3--15). \textbf{Criterion met.}
\item Noise percentage: 15.22\% (SRS target: 5--30\%). \textbf{Criterion met.}
\item Largest cluster share: 44.32\% (SRS target: $\leq 60$\%). \textbf{Criterion met.}
\end{itemize}

\textbf{Approximate spatial cluster inventory:}

\begin{center}
\begin{tabular}{|c|r|l|}
\hline
Cluster & Rows & Corridor assignment \\
\hline
1 & 1{,}097 & Long Island Expressway (I-495) \\
2 & 14{,}792 & Northern State Parkway \\
3 & 11{,}080 & Southern State Parkway \\
4--12 & 90--239 & Secondary corridor segments \\
Noise (0) & 5{,}076 & Unassigned / sparse locations \\
\hline
\end{tabular}
\end{center}

\textbf{Issues:}
\begin{itemize}
\item Corridor assignments (e.g., ``Northern State Parkway'') are determined from
  latitude band ranges (lines 1584--1591). These are heuristic thresholds and
  have not been validated against an authoritative highway shapefile or geocoded
  map. A visual overlay on a base map would substantiate the assignments.
\item Cluster 2 with 14{,}792 rows (44.32\% of all accidents) may represent a
  broad highway band rather than a tightly defined corridor. Splitting this
  cluster with a finer \texttt{eps} could reveal more internal structure.
\item DBSCAN is sensitive to \texttt{eps}. A 0.1 km change in \texttt{eps}
  around the selected value should be tested to confirm robustness of the
  12-cluster solution.
\end{itemize}

\subsection{Phase 8: Post-Clustering Interpretation}

\textbf{Implementation lines:} 1723--2007.

\textbf{What it does:} Joins both cluster label sets back to the full 33{,}375-row
dataset, computes severity, weather, temporal, and roadway profiles by scenario
cluster, computes county composition (Nassau/Suffolk/Queens) by cluster, generates
a spatial-scenario cross-tabulation, and assigns final plain-English names.

\textbf{Verified results:}

\begin{center}
\begin{tabular}{|c|r|r|r|l|}
\hline
Cluster & Rows & Sev3/4 \% & High-risk & Final name \\
\hline
1 & 10{,}389 & 13.99\% & No & Cloudy AM Rush --- Junction Heavy \\
2 & 3{,}857  & 17.27\% & Yes & Scenario Cluster 2 --- Cloudy \\
3 & 10{,}509 & 11.89\% & No & Clear AM Rush --- Junction Heavy \\
4 & 6{,}301  & 16.93\% & Yes & Night / Early AM --- Cloudy (High Risk) \\
5 & 2{,}319  & 18.20\% & Yes & Night / Early AM --- Clear (High Risk) \\
\hline
\end{tabular}
\end{center}

Dataset baseline: Sev3/4 = 12.38\%, Sev4 = 2.75\%.

\textbf{Issues:}
\begin{itemize}
\item Cluster 2 receives a generic name (\textit{Scenario Cluster 2 --- Cloudy})
  because none of the ordered naming conditions in lines 1894--1913 match its
  profile. The naming condition hierarchy should be reviewed to either extend
  coverage or explicitly define a ``generic surface road'' name for
  weather-dominant clusters.
\item Cluster 2 has Sev3/4 = 17.27\%, which is above the dataset baseline of
  12.38\% and above Clusters 1 and 3. The high-risk flag logic (line 1782)
  uses a threshold of 1.5$\times$ baseline Sev4, which is 4.13\%. Cluster 2's
  Sev4 percentage must be below 4.13\% to have been excluded. This threshold
  choice is defensible but should be documented explicitly in the report.
\end{itemize}

\subsection{Phase 9: Export and Visualization}

\textbf{Implementation lines:} 2009--2668.

\textbf{Tangible outputs produced:}

\begin{itemize}
\item \texttt{li\_accidents\_clustered.csv}: 33{,}375 rows $\times$ 58 columns,
  including \texttt{scenario\_cluster}, \texttt{spatial\_cluster},
  \texttt{scenario\_cluster\_name}, \texttt{spatial\_cluster\_name}.
\item \texttt{phase5\_kproto\_fit\_summary.csv}: within-cluster distances and
  diagnostics for $k = 3$--$8$.
\item \texttt{phase6\_validation\_table.csv}: silhouette scores by $k$.
\item \texttt{phase6\_stability\_summary.csv}: Rand and Jaccard indices by $k$.
\item \texttt{phase7\_dbscan\_param\_search.csv}: grid search results across
  all 68 eps $\times$ minPts combinations.
\item \texttt{phase8\_severity\_profile.csv}: Sev1--Sev4 breakdown per cluster.
\item \texttt{phase8\_scenario\_final\_names.csv}: final cluster name table.
\item \texttt{phase8\_spatial\_final\_names.csv}: final spatial cluster name table.
\item Figures: elbow plot, silhouette index plot, stability bar chart,
  risk bubble plot, cluster fingerprint heatmap, spatial hotspot map,
  spatial-scenario cross-tabulation heatmap, county composition bars.
\end{itemize}

\textbf{Issues:}
\begin{itemize}
\item The stability bar chart (Phase 9) uses color zones (green $> 0.75$, amber
  $0.60$--$0.75$, red $< 0.60$) that correctly show both Rand and Jaccard bars
  in the red zone. However, no annotation on the chart or in the surrounding
  report text warns the reader that the red-zone result means the clustering
  should be treated as exploratory.
\end{itemize}

\subsection{Phase 10: Final Report and Submission Checklist}

\textbf{Implementation lines:} 2671--2941.

\textbf{What it does:} Verifies existence of all deliverable files, generates
a limitations table, maps research questions to evidence tables, and runs a
seven-point sanity audit.

\textbf{Sanity audit results:}

\begin{center}
\begin{tabular}{|c|l|l|}
\hline
Check & Criterion & Result \\
\hline
1 & Scenario names unique & \textcolor{passgreen}{PASS} \\
2 & Scenario labels complete (no gaps) & \textcolor{passgreen}{PASS} \\
3 & Spatial labels complete (no gaps) & \textcolor{passgreen}{PASS} \\
4 & Selected Jaccard $\geq 0.4$ & \textcolor{failred}{FAIL} (0.1989) \\
5 & Largest spatial cluster $\leq 60$\% & \textcolor{passgreen}{PASS} (44.32\%) \\
6 & DBSCAN cluster count 3--15 & \textcolor{passgreen}{PASS} (12) \\
7 & DBSCAN noise 5--30\% & \textcolor{passgreen}{PASS} (15.22\%) \\
\hline
\end{tabular}
\end{center}

Check 4 fails with a deficit of $-0.2011$ against the 0.40 threshold. The comment
in the code states that a Jaccard below 0.4 means results should be presented as
exploratory. The limitation table in the same phase does not propagate this
finding into an explicit written statement.

% ============================================================
\section{Critical Defects Requiring Immediate Correction}
% ============================================================

\subsection{Defect 1: Weak Scenario Cluster Quality Not Surfaced to Reader}

\textbf{Severity:} Critical.

\textbf{Location:} Phase 6, Phase 9 stability plot, Phase 10 limitations table.

\textbf{Description:} The scenario clustering solution (k-prototypes, $k=5$)
has silhouette widths of 0.15--0.18 and a Jaccard stability index of 0.1989.
Both values fail standard published thresholds for accepted cluster quality.
The Phase 10 sanity check explicitly detects the Jaccard failure. However:

\begin{itemize}
\item No programmatic warning is printed after Phase 6 completes.
\item The limitations table in Phase 10 attributes uncertainty only to ``sampled
  validation,'' not to the actual metric values observed.
\item The Phase 9 stability bar chart shows red-zone bars but carries no
  explanatory annotation in the chart title or caption.
\item The final report, as currently structured, could present Cluster 4 and
  Cluster 5 as definitive high-risk archetypes without disclosing that the
  clustering framework producing them is statistically weak.
\end{itemize}

\textbf{Required fix:}
\begin{enumerate}
\item Add an explicit printed warning after the Phase 6 selection block:
  \begin{verbatim}
  warning("Scenario clustering solution is exploratory.
           Silhouette = 0.18 (threshold 0.30).
           Jaccard = 0.20 (threshold 0.60).
           Do not present as confirmatory.")
  \end{verbatim}
\item Update the Phase 10 limitations table to include a row stating:
  ``Scenario cluster quality is below standard thresholds
  (silhouette 0.18, Jaccard 0.20). Results are exploratory and should
  not be interpreted as stable accident archetypes.''
\item Add a subtitle or caption note to the Phase 9 stability bar chart
  indicating that both metrics fall in the weak zone.
\end{enumerate}

\subsection{Defect 2: k-Selection Fallback Logic Produces Incorrect k}

\textbf{Severity:} High.

\textbf{Location:} Phase 6, approximately line 1288.

\textbf{Description:} The selection logic reads:
\begin{verbatim}
if (length(selected_stability) == 1 &&
    selected_stability < 0.5 &&
    nrow(stability_summary) > 1) {
  selected_k <- stability_summary[order(-jaccard, -rand, k)][1, k]
}
\end{verbatim}

With \texttt{selected\_stability = 0.1989}, all three conditions should be
satisfied and the fallback to $k=7$ (higher Jaccard = 0.1666 is actually
\textit{lower} than $k=5$'s 0.1989, so the fallback would select $k=5$ on
Jaccard ordering anyway --- but the logic path itself needs verification).

On closer inspection: when \texttt{stability\_summary} is ordered by
\texttt{-jaccard}, $k=5$ (Jaccard 0.1989) ranks first, so the fallback would
still select $k=5$. The logic may be working correctly in its own terms, but:

\begin{itemize}
\item The fallback only tests two $k$ values (5 and 7) because \texttt{local\_safe}
  profile runs bootstrap only on candidates near the silhouette-best.
\item The design intent (fall back to highest Jaccard) and the actual outcome
  (stays at $k=5$ because it has higher Jaccard than $k=7$) are consistent.
  However, both candidates have Jaccard $< 0.2$, so choosing between them is
  analytically meaningless.
\end{itemize}

\textbf{Required fix:}
\begin{enumerate}
\item Add a diagnostic print block after line 1289 that explicitly logs:
  which $k$ values were evaluated for stability, their Jaccard values, which
  condition triggered or did not trigger, and the final selected $k$.
\item Add an explicit check: if \texttt{max(stability\_summary\$jaccard) < 0.4},
  print a named warning identifying the best available $k$ as still failing
  the threshold.
\end{enumerate}

\subsection{Defect 3: run\_profile Set to local\_safe for Final Submission}

\textbf{Severity:} High.

\textbf{Location:} Phase 1, line 135.

\textbf{Description:} The current profile (\texttt{local\_safe}) uses a 5{,}000-row
sample (15\% of data) for validation with a single bootstrap repeat. Silhouette
and Jaccard estimates derived from 15\% of the data are noisy and unreliable as
summary statistics of the full 33{,}375-row dataset. The same profile that
produced the failing Jaccard of 0.1989 may show a different pattern on the full
dataset. The observed quality failure could be an artifact of under-sampling.

\textbf{Required fix:}
\begin{enumerate}
\item Before final report submission, rerun with \texttt{run\_profile <- "consensus\_16gb"}.
  This uses 4{,}000-row samples averaged over 4 repeats and 10 bootstrap runs.
\item Document the updated silhouette and Jaccard values in the report. If the
  quality failure persists under \texttt{consensus\_16gb}, it must be disclosed
  prominently. If quality improves, the current \texttt{local\_safe} values should
  be footnoted as preliminary.
\item If hardware memory allows (available RAM $> 16$ GB and R vector limit
  raised), consider running \texttt{full\_srs} for at least the validation phase.
  Note: the full-dataset silhouette computation hit the 24 GB R memory limit
  in testing and requires either a higher \texttt{mem.maxVSize()} setting or
  a subsampled silhouette approximation.
\end{enumerate}

% ============================================================
\section{Improvements Recommended Before Final Report}
% ============================================================

\subsection{High Priority}

\textbf{IP-1: Feature Set Sensitivity Analysis.}

The scenario clustering uses 16 features including \texttt{Pressure(in)} and
\texttt{Visibility(mi)}, both of which required median imputation. Including
imputed values in a distance-based clustering method can deflate separation
by pulling imputed points toward the global median artificially.

\textbf{Action:} Rerun Phase 5 with a reduced feature set that excludes
\texttt{Pressure(in)} and \texttt{Visibility(mi)}. Compare silhouette and
Jaccard outcomes. If the reduced feature set produces meaningfully higher
scores, adopt it and document the decision.

\textbf{IP-2: Verify DBSCAN Corridor Assignments Against Geography.}

The current corridor labels (e.g., ``Northern State Parkway Cluster 2'') are
assigned by latitude band thresholds coded at lines 1584--1591. These are
educated guesses based on approximate highway latitude values. A misassignment
would be embarrassing in a final report.

\textbf{Action:} Print the centroid coordinates of each spatial cluster in
km and degrees. Manually verify in a mapping tool (Google Maps or similar)
that each centroid falls on or near the named highway. Update the label
if any assignment is wrong.

\textbf{IP-3: Add Confidence Interval to Jaccard Estimates.}

The \texttt{local\_safe} profile runs one bootstrap repeat. A single repeat
gives a point estimate of Jaccard with no variance measure. The same run
can produce meaningfully different Jaccard values if the random sample is
different.

\textbf{Action:} Under \texttt{consensus\_16gb} (4 repeats) or higher, compute
the standard deviation of Jaccard across bootstrap runs and report as
$\bar{J} \pm \sigma_J$.

\subsection{Medium Priority}

\textbf{IP-4: Cluster 2 Naming Improvement.}

Cluster 2 (``Scenario Cluster 2 --- Cloudy'') is the only cluster without a
meaningful scenario-level name. The Phase 8 naming hierarchy exhausts
conditions for Winter, Night, Rush, and Weekend profiles but never assigns
a surface-road midday pattern. Cluster 2 has 3{,}857 rows and Sev3/4 of
17.27\%, making it the second-largest risk cluster.

\textbf{Action:} Inspect the dominant hour\_band and Sunrise\_Sunset values
for Cluster 2. Add a condition to the naming hierarchy that captures
``Cloudy Midday Surface Road'' or similar based on those values.

\textbf{IP-5: Report Severity 4 Rate by Cluster Explicitly.}

The current risk profiling uses Sev3/4 combined as the risk metric.
Severity 4 (the most severe category) represents 2.75\% of the dataset.
Clusters 4 and 5 are flagged as high-risk, but the magnitude of the
Severity 4 elevation over baseline is not stated explicitly in a table.

\textbf{Action:} Add a column to the severity profile table reporting the
Severity 4 rate per cluster alongside the Sev3/4 rate. Include the
dataset baseline (2.75\%) as a reference row.

\textbf{IP-6: eps Sensitivity Check for DBSCAN.}

The selected \texttt{eps} = 1.0 km falls in the middle of the tested grid.
The 12-cluster result may change significantly with \texttt{eps} = 0.9 km
or 1.1 km.

\textbf{Action:} Run two additional DBSCAN trials at $\text{eps} = 0.9$ and
$\text{eps} = 1.1$ (holding \texttt{minPts} = 100). Report whether the
cluster count and corridor structure are stable, or whether the solution
is sensitive to the eps value.

\subsection{Low Priority}

\textbf{IP-7: Binary Flag Factor Labeling (Line 585--589).}

The retained binary flags are converted to factor using
\texttt{factor(x, labels = c("FALSE", "TRUE"))}. This is unconventional.
The standard idiom is \texttt{as.factor()} on a logical column, which
automatically assigns levels \texttt{FALSE} and \texttt{TRUE}. The current
approach works but is misleading to a reader reviewing the code.

\textbf{Action:} Replace with \texttt{as.factor()} or
\texttt{factor(x, levels = c(FALSE, TRUE))}.

\textbf{IP-8: Thunderstorm Weather Group Dependency.}

The weather grouping for Thunderstorm (line $\approx$ 533) relies on
\texttt{grepl("Thunderstorm", x, ignore.case = TRUE)} rather than an explicit
entry in the lookup table used for other groups. This means a value like
``Thunder Storm'' (with a space) would fall into the \texttt{Other} category.

\textbf{Action:} Add common Thunderstorm string variants to the explicit lookup
map rather than relying solely on \texttt{grepl()}.

% ============================================================
\section{What Works and What Does Not}
% ============================================================

\subsection{What Works}

\begin{itemize}
\item \textbf{DBSCAN spatial clustering} produces a credible and well-validated
  result: 12 corridor clusters, 15.2\% noise, all SRS targets met. The highway
  corridor structure of Long Island is real and the algorithm has detected it
  correctly.
\item \textbf{Feature engineering} is thorough. Temporal features, weather
  grouping, wind normalization, and binary flag filtering are all methodologically
  sound.
\item \textbf{Data integrity checks} are rigorous. Dimension assertions, missingness
  profiling, and post-imputation NA validation make the pipeline robust against
  silent data corruption.
\item \textbf{Reproducibility} is solid. Seeded, profiled, and parametrically
  documented runs can be replicated exactly.
\item \textbf{Post-clustering interpretation} (Phase 8) is valuable and correct.
  The severity profiles and county distributions are real findings that stand
  independently of the cluster quality concern.
\item \textbf{Sanity audit} (Phase 10) is an honest self-check that correctly
  identifies the Jaccard failure. This is methodological integrity.
\end{itemize}

\subsection{What Does Not Work}

\begin{itemize}
\item \textbf{Scenario clustering quality} does not meet standard thresholds.
  The k-prototypes solution has weak cluster separation (silhouette $\approx$ 0.18)
  and poor stability (Jaccard $\approx$ 0.20). The result is a real partition
  of the data, but it is not a \textit{stable} partition: different random samples
  of the data do not reliably reproduce the same five clusters.
\item \textbf{The cluster quality failure is not communicated} to the reader
  of the rendered report. Only the code comments acknowledge it.
\item \textbf{The validation was run on a small sample} (5{,}000 rows,
  \texttt{local\_safe} profile). The final silhouette and Jaccard values may
  not accurately represent the full dataset.
\end{itemize}

% ============================================================
\section{Bottom Line and Submission Readiness}
% ============================================================

\subsection{Spatial Clustering Track: Ready with Minor Improvements}

The DBSCAN component meets all SRS acceptance criteria. The spatial findings ---
12 highway corridor hotspots along the LIE, Southern State Parkway, and Northern
State Parkway --- are defensible and original. Before submission, corridor
assignments should be manually verified against geographic coordinates
(Improvement IP-2) and an eps sensitivity check should be documented (IP-6).

\subsection{Scenario Clustering Track: Not Submission-Ready Without Disclosure}

The k-prototypes component produces a $k=5$ partition that is real but weak.
The following must be done before the scenario clustering results can appear
in a final academic report:

\begin{enumerate}
\item Rerun with \texttt{consensus\_16gb} profile to get more reliable
  validation estimates (Defect 3).
\item If quality metrics improve: report the updated values and note that
  the earlier \texttt{local\_safe} run was preliminary.
\item If quality metrics remain poor: add an explicit exploratory disclosure
  block to the report (Defect 1). Do not present Clusters 4 and 5 as
  proven high-risk archetypes; present them as candidate high-risk patterns
  that warrant further investigation.
\item Fix the k-selection diagnostic logging (Defect 2).
\item Attempt the feature set reduction (IP-1) to test whether removing
  imputed features tightens the cluster structure.
\end{enumerate}

\subsection{Overall Academic Defensibility Rating}

\begin{center}
\begin{tabular}{|l|l|}
\hline
\textbf{Component} & \textbf{Defensibility} \\
\hline
Spatial (DBSCAN) results & \textcolor{passgreen}{Strong --- ready for report} \\
Feature engineering & \textcolor{passgreen}{Sound --- no changes needed} \\
Post-clustering profiles & \textcolor{passgreen}{Valid --- useful findings} \\
Scenario clustering results & \textcolor{warnamber}{Exploratory only --- requires disclosure} \\
Validation under local\_safe & \textcolor{failred}{Insufficient --- must rerun at higher profile} \\
\hline
\end{tabular}
\end{center}

% ============================================================
\section{References}
% ============================================================

\begin{thebibliography}{9}

\bibitem{kaufman1990}
Leonard Kaufman and Peter J.\ Rousseeuw.
\textit{Finding Groups in Data: An Introduction to Cluster Analysis}.
Wiley, 1990. Defines the silhouette width interpretation thresholds (0.25, 0.50, 0.70).

\bibitem{hennig2007}
Christian Hennig.
\textit{Cluster-wise assessment of cluster stability}.
Computational Statistics and Data Analysis, 52(1):258--271, 2007.
Defines the Jaccard stability thresholds (0.60 acceptable, 0.75 stable).

\bibitem{szepannek2018}
Gero Szepannek.
\textit{clustMixType: User-Friendly Clustering of Mixed-Type Data in R}.
The R Journal, 2018.
\url{https://journal.r-project.org/articles/RJ-2018-048/}

\bibitem{hahsler2019}
Michael Hahsler, Matthew Piekenbrock, and Derek Doran.
\textit{dbscan: Fast Density-Based Clustering with R}.
Journal of Statistical Software, 91(1), 2019.
\url{https://www.jstatsoft.org/v91/i01/}

\bibitem{r-validation}
R Documentation.
\textit{clustMixType::validation\_kproto}.
\url{https://search.r-project.org/CRAN/refmans/clustMixType/html/validation_kproto.html}

\bibitem{r-stability}
R Documentation.
\textit{clustMixType::stability\_kproto}.
\url{https://search.r-project.org/CRAN/refmans/clustMixType/html/stability_kproto.html}

\end{thebibliography}

\end{document}
